% The finite Dirac operator as a boundary condition: short note, from lemmas 22 and 23
\ifdefined\mobile
\documentclass[aps,prd,reprint,onecolumn,nofootinbib,amsmath,amssymb,floatfix]{revtex4-2}
\usepackage[paperwidth=270pt,paperheight=585pt,textwidth=246pt,textheight=545pt,top=18pt,left=12pt,headheight=0pt,headsep=0pt,footskip=14pt]{geometry}
\else
\documentclass[aps,prd,reprint,nofootinbib,amsmath,amssymb,floatfix]{revtex4-2}
\fi
\usepackage[T1]{fontenc}
\usepackage{lmodern}
\usepackage{amsthm}
\emergencystretch=1.5em
\ifdefined\mobile\makeatletter\def\frontmatter@abstractwidth{246\p@}\makeatother\fi
\usepackage[hidelinks]{hyperref}
\theoremstyle{plain}
\newtheorem{theorem}{Theorem}
\newtheorem{proposition}[theorem]{Proposition}
\newtheorem{corollary}[theorem]{Corollary}
\newcommand{\C}{\mathbb C}
\renewcommand{\H}{\mathbb H}
\newcommand{\supp}[1]{Supplemental Material, part #1}

\begin{document}

\title{The finite Dirac operator as a boundary condition on an interval}
\date{Draft, 3 November 2026}

\begin{abstract}
Every finite Dirac operator $D_F$ with the signs of KO-dimension six is exactly the lowest spectral window of the pure derivative on an interval of length $l$, for $l\|D_F\|<\pi/2$, with no mass term. The boundary data are a grading at each wall, and $D_F=(2l)^{-1}\Gamma_0\log(\Gamma_1\Gamma_0)$, so the singular values of $D_F$ are principal angles over $l$ and the masses are quantized momenta. Preserving the domain restricts the observables at a wall to the commutant of its grading. Order one holds in this parent for any finite data with order zero. The compression of an observable is its average under $e^{ixD_F}$, $|x|\le l$, which tends to the finite representation with rate $l^2$. Aligned walls compress a bulk prealgebra to its even part, with the grading inherited. Even scalar fluctuations are pure gauge, and on a moving window they reproduce the quadratic term of Chamseddine-Connes-van Suijlekom. Odd ones change the angles; for $M_2(\H)$ their field space is $S^4$. In a Yang-Mills parent with a wall connector the $W$ and $Z$ masses are exact, with the constants $2$ and $16/5$ at small angle. All finite data are inputs and the parent algebra is matrix-valued, so the selection of the finite algebra is untouched.
\end{abstract}

\maketitle

\section{Introduction}

In an almost-commutative geometry the masses of the fermions are the entries of a finite Dirac operator $D_F$, and in a Kaluza-Klein reduction a mass can be a momentum in the internal direction, quantized by the boundary conditions. We relate the two for an interval. The Connes-Lott mass is the lowest such momentum, $\pi/(2l)$, and the known boundary formula $\tan(m_nL)=M_DL$ of Higgsless models \cite{Csaki2004} gives the interpolation for one flavor. Here we carry this through for the whole finite operator, with its real structure and grading, and we ask what the parent says about the algebra, order one, and fluctuations.

The finite data remain inputs, so the existence of the window is close to a reparametrization. The content lies in the existence of a consistent parent, and in the form it gives to the wall algebra, the compression and the fluctuations.

\section{The window}

Let $H_F=H_+\oplus H_-$ be finite-dimensional, $\Gamma_F=\mathrm{diag}(1,-1)$, and $c:H_+\to H_-$ antiunitary, with $J_F(a,b)=(c^{-1}b,ca)$. Let $D_F$ have blocks $M$ and $M^*$, with $M=cM^*c$, which is $J_FD_F=D_FJ_F$. These are the signs $(+,+,-)$. For $l\|D_F\|<\pi/2$, define $X:H_-\to H_+$ as the block of $\tan(lD_F)$, so that $X^*$ is the other block.

On $L^2([0,l];\C^2)\otimes(H_+\oplus H_-)$, write $\psi$ with components $(Q_L,Q_R;U_L,U_R)$, and let $D_X=E\,d/dx$, where $(E\psi)_L=\psi_R$ and $(E\psi)_R=-\psi_L$, with domain
\begin{equation}
\begin{aligned}
&Q_R(0)=U_L(0)=0,\\
&Q_R(l)=XU_R(l),\quad U_L(l)=-X^*Q_L(l).
\end{aligned}
\end{equation}
Let $\hat\Gamma$ be $+1$ on $\psi_L$ and $-1$ on $\psi_R$, and let $\hat J\psi$ have components $(c^{-1}U_R,-c^{-1}U_L;-cQ_R,cQ_L)$.

\begin{theorem}\label{thm:window}
$D_X$ is self-adjoint with compact resolvent, with signs $(+,+,-)$. Let $R(x)=\exp(x\Gamma_FD_F)$, and let $W:H_F\to L^2$ send $(a,b)$ to the function with $L$-components $l^{-1/2}R(x)(a,0)$ and $R$-components $l^{-1/2}R(x)(0,b)$. Then $W$ is an isometry, and
\begin{equation}
D_XW=WD_F,\quad \hat JW=WJ_F,\quad \hat\Gamma W=W\Gamma_F.
\end{equation}
With $\theta_k=l\sigma_k$ for the singular values $\sigma_k$ of $M$, the signed spectrum is $(n\pi\pm\theta_k)/l$, $n\in\mathbb Z$, so $\mathrm{ran}\,W$ is the spectral subspace of $|D_X|\le\|D_F\|$, and the next level is at least $(\pi-l\|D_F\|)/l$.
\end{theorem}

\begin{proof}[Structure of the proof]
Takagi factorization gives $M=V|M|$ with $V=cV^*c$. Then $R(x)(a,0)$ has components $\cos(x|M|)a$ and $-V\sin(x|M|)a$, and $R(x)(0,b)$ has components $\sin(x|M|)V^*b$ and $\cos(x|M^*|)b$. The wall conditions follow from $X=\tan(l|M|)V^*$, and $D_XW=WD_F$ follows by differentiating $R$. Constant unitaries on $Q$ and $U$ reduce $X$ to $\mathrm{diag}(\tan\theta_k)$, and each channel has the known one-flavor spectrum (\supp{22}).
\end{proof}

So $D_X$ has no mass term. A light eigenvalue is a momentum $k=\theta/l$, fixed by $\tan(kl)=\tan\theta$, a singular value of $X$. In the finite description this condition is lost and the value appears as free data, which offers a reading of why the finite triple needs Yukawa couplings. One matrix $X$ encodes the Yukawa and Majorana blocks together. Also, $R^*D_XR=D_0+E\Gamma_FD_F$, with $D_0$ uncoupled: after a unitary change of domain, $D_F$ is a constant matrix connection, with relative endpoint rotation $R(l)$.

\begin{theorem}\label{thm:gradings}
The boundary conditions are gradings: $\psi_L$ lies in the $+1$ eigenspace of $\Gamma_i$ and $\psi_R$ in the $-1$ eigenspace at each wall, with $\Gamma_0=\Gamma_F$ and $\Gamma_1$ equal to $R(l)\Gamma_0R(l)^*$. If $-1$ is not an eigenvalue of $\Gamma_1\Gamma_0$, then
\begin{equation}
D_F=\frac1{2l}\,\Gamma_0\log(\Gamma_1\Gamma_0),
\end{equation}
and $l\sigma_k$ are the principal angles between the $+1$ eigenspaces. A real structure is given by an antiunitary $J_1$, with $J_1^2=-1$, that anticommutes with $\Gamma_0$ and $\Gamma_1$. Then $J_F=\Gamma_0J_1$, and $J_F$ carries $\Gamma_1$ to $-\Gamma_0\Gamma_1\Gamma_0$.
\end{theorem}

At an angle $\pi/2$ the channel has the chiral conditions of the Connes-Lott interval, with no zero mode and lowest level $\pi/(2l)$.

\section{Algebra, order one, compression}

Let a unital algebra $A_F$ act evenly and faithfully, with order zero, and let $C_F$ be the commutant of $D_F$ in $A_F$. Order one is not assumed.

\begin{theorem}\label{thm:parent}
A function $f\in C^\infty([0,l],A_F)$ preserves the domain of $D_X$ iff $f(l)\in C_F$. Let $A_l$ be the algebra of such $f$. Then $[D_X,\pi(f)]=E\pi(f')$, and we have a real even spectral triple with order zero and order one.
\end{theorem}

\begin{proof}[Structure of the proof]
The wall conditions say that $\pi(f(l))$ commutes with $\tan(lD_F)$, equivalently with $D_F$, for $l\|D_F\|<\pi/2$. The opposite action is pointwise, so by finite order zero $E\pi(f')$ commutes with it.
\end{proof}

So $A_F$ acts freely in the bulk, and the wall carries the commutant of $D_F$. In the Standard Model data with nonzero Yukawa blocks, this is $\{(\lambda,\lambda,m)\}$, the algebra of the photon and gluons. Finite order one does not arise as a condition on the parent.

\begin{theorem}\label{thm:compression}
Let $Q_l(T)=W^*TW$. For even $T(x)$,
\begin{equation}
\begin{aligned}
Q_l(T)&=\frac1l\int_0^l\cos(x\,\mathrm{ad}_{D_F})(T(x))\,dx,\\
Q_l(ET)&=-\frac1l\int_0^l\sin(x\,\mathrm{ad}_{D_F})(T(x))\,dx.
\end{aligned}
\end{equation}
For constant $a$,
\begin{equation}
\begin{aligned}
Q_l(a)&=\frac1{2l}\int_{-l}^{l}e^{ixD_F}\pi(a)e^{-ixD_F}dx\\
&=\mathrm{sinc}(l\,\mathrm{ad}_{D_F})(\pi(a)),
\end{aligned}
\end{equation}
so $\|Q_l(a)-\pi(a)\|\le(l^2/6)\|[D_F,[D_F,\pi(a)]]\|$. Moreover $Q_l(A_l)=Q_l(C^\infty([0,l],A_F))$, which contains $Q_l(a)$ for every $a\in A_F$, and the order defects of the compressed data differ from the finite ones by $O(l^2)$ with explicit constants.
\end{theorem}

Thus $D_F,J_F,\Gamma_F$ are exact, and the observables converge to $\pi(A_F)$ as $l\to0$ at fixed $D_F$, as maps. The generated algebra is discontinuous at $l=0$: compressions of commuting projections can fail to commute at order $l^2$, in line with the remark of Connes and van Suijlekom \cite{CvS2021} on generated algebras (\supp{22}).

\begin{theorem}\label{thm:aligned}
Let $P$ be a $*$-subalgebra of $\mathrm{End}(H)$ that need not commute with $\Gamma_0$, acting pointwise and identically on both spin components. A function $f$ valued in $P$ preserves the domain iff $[f(0),\Gamma_0]=0$ and $[f(l),\Gamma_1]=0$, and $\hat\Gamma$ commutes with every such $f$. For aligned walls $\Gamma_1=\Gamma_0$, $D_F=0$, $\hat\Gamma$ restricts to $\Gamma_0$ on the window, and $Q(a)=(a+\Gamma_0a\Gamma_0)/2$.
\end{theorem}

So a bulk prealgebra is compressed to its even part, with the grading inherited from the parent operator $\hat\Gamma$. For $M_2(\H)\oplus M_4(\C)$ this gives $\H\oplus\H\oplus M_4(\C)$, the step of the classification \cite{CC2008} that a compression of functions on $S^4$ reaches only with an added grading.

\section{Fluctuations}

\begin{theorem}\label{thm:transport}
Let $h$ be smooth with anti-Hermitian values in $P$, with $u'=-hu$, $u(0)=1$, $w=u(l)$, and let $D(\Gamma_0,\Gamma_1)$ be the derivative with those gradings. Then $D(\Gamma_0,\Gamma_1)+Eh$ equals $uD(\Gamma_0,w^*\Gamma_1w)u^*$. The self-adjoint scalar one-forms are exactly the $Eh$. If $h$ takes values in the even algebra, then $w^*\Gamma_1w$ has the same angles with $\Gamma_0$; on the moving window the finite operator is $\pi(w)^*D_F\pi(w)$, and with the real fluctuation $U_w^*D_FU_w$, where $U_w=\pi(w)J_F\pi(w)J_F^{-1}$. In that case the connection is gauge equivalent to zero under the allowed endpoint gauge group.
\end{theorem}

With $v=w^*$, this is the gauge orbit of \cite{CCvS2013a}, Eqs.~(10)-(11), namely $UDU^*$, which equals $D+A_{(1)}+\tilde A_{(1)}+A_{(2)}$, quadratic term included, for finite data that fail order one, while the parent has order one and no quadratic term. The fixed compression is a different object: for constant $h$, $Q_l(Eh)=-(l/2)[D_F,h]+O(l^3)$, against $-l[D_F,h]$ for the first derivative of the transported operator.

\emph{No-go theorem (the radial mode).} For one interval, fixed $X$, no bulk masses, even scalar inner fluctuations and the moving window, the singular values of the finite operator are fixed. With nonzero Yukawa data, the Higgs quaternion stays on the unit quaternions, with no radial mode (\supp{22}). The escapes are the following: $X$ as a field on the wall; bulk terms odd for the grading; bulk masses not commuting with the gauge action; more internal dimensions; an added composite degree of freedom. At the $10\to9$ endpoint this agrees with a frozen Higgs field.

The odd escape is explicit. For $W=\H^2$, the gradings with quaternionic eigenspaces form $S^4=Sp(2)/(Sp(1)\times Sp(1))$, with $\Gamma_p=p_a\Gamma_a$. For two points at geodesic distance $d$, all principal angles equal $\theta=d/2$. On $\mathrm{Hom}(V,W)\oplus\mathrm{Hom}(W,V)$, the connection $K(x,y)=(\theta/l)(Fx,yF)$, with $F$ odd in $M_2(\H)$ and $F^2=-1$, gives a universal mass $\theta/l$, with $J_1=\Gamma_0J$ and a second grading at the point $p$ on the first summand and its mirror image on the second. So $(\theta,q)$ are polar coordinates of a Higgs quaternion on $S^4$, the coset of the minimal composite Higgs model on an interval \cite{ACP2005}, in the form of gauge-Higgs unification \cite{Hosotani1983}. That $S^4$ is also the base of the twistor window is an observation here, with no identification. $J_1$ induces on odd elements an opposite action twisted by the grading automorphism, which is the automorphism of the twisted triples of the grand symmetry \cite{DLM2014}; the comparison of their first-order conditions is open.

\section{Gauge sector and the two limits}

Assume a flat product and a positive five-dimensional Yang-Mills kinetic form with no boundary terms. Then $A_\mu$ at a wall lies in the wall algebra, the directions in it are Neumann and the others Dirichlet, as in Higgsless models \cite{Csaki2003}.

\emph{The strict family.} For even walls with nonzero Yukawa data, the unbroken directions have masses $n\pi/l$ and the broken ones $(n-\tfrac12)\pi/l$, so $M_W=M_Z=\pi/(2l)$, with massless photon and gluons after unimodularity. Fix a light block $M_0$ and a heavy block $\Theta/l$ with $0<\Theta<\pi/2$. As $l\to0$, the light couplings $l|M_0|$ vanish at fixed mass, the heavy mass $\Theta/l$ and $M_W$ diverge with $m_t/M_W=2\Theta/\pi<1$, and the light compression converges. This realizes the $10\to9$ proposal at this level, with $X_q\sim(\pi/2)m_q/M_W$, in the form of $y_q=\sqrt2m_q/v\to0$ at fixed mass. Identifying $X_q$ with a four-dimensional Yukawa coupling is a hypothesis.

\begin{theorem}\label{thm:wz}
Take $P=M_2(\H)\oplus M_4(\C)$, with $sp(2)\oplus su(4)$ and the trace form, with one coupling, gradings at distance $2\theta$, and at $x=l$ a connector on a line $w$ in the $+1$ eigenspace of $\Gamma_p$ and a line $v$ in $\C^4$, whose stabilizer is $\C\oplus\H\oplus M_3(\C)$. Then exactly nine modes are massless, the photon, with $1/e^2=(8/3)/g^2$, and the gluons, and
\begin{equation}
\begin{aligned}
\tan^2(M_Wl)&=\frac{\sin^22\theta}{1+\cos^22\theta},\\
\tan^2(M_Zl)&=\frac{4\sin^22\theta}{1+4\cos^22\theta}.
\end{aligned}
\end{equation}
With $m=\theta/l$, as $\theta\to0$: $M_W^2\to2m^2$ and $M_Z^2\to\tfrac{16}5m^2$. As $\theta\to\pi/2$: $M_W,M_Z\to0$ and $m\to\pi/(2l)$. In both limits $M_W^2/M_Z^2\to5/8$.
\end{theorem}

The constants $2$ and $16/5$ are $\|[D_F,a]\|^2/\|a\|^2$, and they are the lower ends of the poles of the fully massive family in the exact quadratic spectral action (\supp{18}), whose Taylor mass carries the additional factor $3/2$ from the heat coefficients. The second limit is a geometric form of the behavior found there, with heavy fermions and light vectors: at antipodal points the two walls have the same stabilizer in $sp(2)$.

\begin{proposition}\label{prop:theta}
For the heat cutoff and universal $\theta$, with $s=l\Lambda$, the trace per channel is $2\sum_{n\in\mathbb Z}\exp(-(n\pi+\theta)^2/s^2)$, a Jacobi theta function, strictly decreasing on $[0,\pi/2]$. For $s\to0$ at fixed $x=\theta^2/s^2$, it tends to $2e^{-x}$.
\end{proposition}

So the runaway of the radial direction in the finite action \cite{towerlesspaper} has an escape on the interval: the radial direction is compact, and the tower stops it at the stationary point $m=\pi/(2l)$. There, by Theorem~\ref{thm:wz}, the vectors are massless, so this is not the physical point. In the same family, a numerical scan of the one-loop potential with vector and fermion loops, including bulk masses for a top channel, finds no interior minimum (Python, \supp{23}). In a flat tree-level comparison, neither the Wilson-line family, with $M_W^2/M_Z^2\in[0.503,0.625]$, nor a finite wall scalar, with $[0.625,1]$, fits both the top ratio and the weak ratio. These negative statements hold in those flat families. They point toward a wall scalar with its own potential, or warping, for the physical interior, consistently with the reading of the odd labels as limits.

\section{Status}

As mathematics, the one-flavor formula and the gauge boundary conditions are known \cite{Csaki2004,Csaki2003}, compression of operator systems is treated in \cite{CvS2021}, and spectral action on spaces with boundary in \cite{CC2011}, for a boundary of the manifold, with a local chirality projector. In our searches we did not find the finite real triple as an exact window on an internal interval, with the wall algebra and the average formula; the searches are recorded in the Supplemental Material and do not establish absence. The parent algebra is matrix-valued, so the founding question for a commutative parent remains open, as does any selection of $A_F$.

\section*{Provenance}

This note has no listed authors. The results were proved by Claude Opus 5.5, and refereed by Terra 6.1 (gpt-6.1-terra, effort max), whose reports led to corrections, with Python (numpy, scipy) checks of the identities and the numerical statements. Details are in the Supplemental Material of \cite{towerlesspaper}, parts 22 and 23, with the referee reports.

\bibliography{towerless_intervals}

\end{document}
